\documentclass[11pt,a4paper]{article}
\usepackage[utf8]{inputenc}
\usepackage[T1]{fontenc}
\usepackage[french]{babel}
\usepackage{amsmath, amssymb}
\usepackage{geometry}
\geometry{margin=2cm}
\usepackage{tabularx}
\usepackage{booktabs}
\usepackage{xcolor}
\usepackage{prftree}
\usepackage{graphicx}

\title{\vspace{-2cm}\textbf{Cheatsheet Rocq : Leçon 1}}
\date{}

\begin{document}
\maketitle
\vspace{-1.5cm}

\section*{1. Dictionnaire : $\lambda$-calcul simplement typé $\leftrightarrow$ Rocq}

\renewcommand{\arraystretch}{1.5}
\begin{tabularx}{\textwidth}{@{} l X c @{}}
  \toprule
  \textbf{Concept Théorique} & \textbf{Notation Papier (Maths/Logique)} & \textbf{Syntaxe Rocq (Gallina)} \\
  \midrule
  %   Type des propositions & --- & \texttt{Prop} \\
  Déclaration de type & $x : A$ & \texttt{x : A} \\
  Implication / Type flèche & $A \to B$ & \texttt{A -> B} \\
  Abstraction ($\lambda$) & $\lambda x. M$ & \texttt{fun x => M} \\
  Abstraction multiple & $\lambda x. \lambda y. M$ & \texttt{fun x y => M} \\
  Application & $M \ N$ & \texttt{M N} \\
  Parenthésage & $(f \ (g \ x))$ & \texttt{f (g x)} \\
  \bottomrule
\end{tabularx}

\vspace{1cm}

\section*{2. Tactiques et Déduction Naturelle}

Dans le mode \textbf{Preuve Interactive}, les tactiques construisent le $\lambda$-terme à l'envers (en remontant de la conclusion vers les prémisses de l'arbre de dérivation).

\vspace{0.5cm}
\begin{tabularx}{\textwidth}{@{} l X X @{}}
  \toprule
  \textbf{Tactique Rocq} & \textbf{Effet sur le Séquent (Contexte $\vdash$ But)} & \textbf{Règle de Déduction Naturelle} \\
  \midrule
  \texttt{intros x.} &
  Prend l'antécédant $A$ du but ($A \to B$) et l'introduit comme nouvelle hypothèse dans le contexte sous le nom \texttt{x}. Le but devient $B$. &
  \textbf{$\to$-Introduction} \newline (Abstraction en $\lambda$-calcul) \newline
  \vspace{0.2cm} \newline
  $\prftree[r]{Abs}{\Gamma, x : A \vdash M : B}{\Gamma \vdash \lambda x.M : A \to B}$ \\
  \midrule
  \texttt{exact x.} &
  Résout le but courant si l'hypothèse \texttt{x} correspond \emph{exactement} au but. &
  \textbf{Règle "axiome"} \newline (Typage d'une variable) \newline
  \vspace{0.2cm} \newline
  $\prftree[r]{Var}{\Gamma, x : A \vdash x : A}$ \\
  \midrule
  \texttt{apply f.} &
  Si le but est $B$ et qu'on a $f : A \to B$, change le but en $A$ (raisonnement arrière). &
  \textbf{$\to$-Élimination} \newline (Modus Ponens / Application) \newline
  \vspace{0.2cm} \newline
  \resizebox{\linewidth}{!}{$\prftree[r]{App}{\prftree[r]{Var}{\Gamma, f : A \to B \vdash f : A \to B}}{\prfassumption{\Gamma, f : A \to B \vdash M : A}}{\Gamma, f : A \to B \vdash f\ M : B}$} \\
  \bottomrule
\end{tabularx}

\end{document}